% ============================================================================
% main.tex — 《模形式与费马大定理》
% ============================================================================
\documentclass[12pt, a4paper, openany]{book}

% ============================================================================
% preamble.tex — 包、字体、排版、数学环境
% ============================================================================

% --- 中文支持 (XeLaTeX) ---
\usepackage{ctex}              % 中文排版核心
\usepackage{fontspec}          % 字体选择

% --- 页面布局 ---
\usepackage[
  a4paper,
  top=2.5cm, bottom=2.5cm,
  left=2.5cm, right=2.5cm
]{geometry}

% --- 数学 ---
\usepackage{amsmath, amssymb, amsthm}
\usepackage{mathrsfs}          % \mathscr 花体
\usepackage{mathtools}         % dcases, coloneqq 等
\usepackage{bm}                % 粗体数学符号

% --- 图表 ---
\usepackage{graphicx}
\usepackage{float}
\usepackage{booktabs}
\usepackage{longtable}
\usepackage{tabularx}

% --- 绘图 ---
\usepackage{tikz}
\usetikzlibrary{
  positioning, calc, arrows.meta,
  decorations.pathreplacing,
  patterns,
}

% --- 交叉引用与超链接 ---
\usepackage{hyperref}
\hypersetup{
  colorlinks=true,
  linkcolor=blue!60!black,
  citecolor=green!40!black,
  urlcolor=cyan!50!black,
}
\usepackage[nameinlink]{cleveref}  % 智能引用: \cref

% --- 索引与参考文献 ---
\usepackage{makeidx}
\makeindex

% --- 颜色 ---
\usepackage{xcolor}

% ============================================================================
% 定理环境
% ============================================================================
\theoremstyle{plain}
\newtheorem{theorem}{定理}[chapter]
\newtheorem{proposition}[theorem]{命题}
\newtheorem{lemma}[theorem]{引理}
\newtheorem{corollary}[theorem]{推论}
\newtheorem{conjecture}[theorem]{猜想}

\theoremstyle{definition}
\newtheorem{definition}[theorem]{定义}
\newtheorem{example}[theorem]{例}
\newtheorem{exercise}{习题}[chapter]

\theoremstyle{remark}
\newtheorem{remark}[theorem]{注记}
\newtheorem{notation}[theorem]{记号}

% ============================================================================
% 常用数学命令
% ============================================================================

% 数域
\newcommand{\N}{\mathbb{N}}
\newcommand{\Z}{\mathbb{Z}}
\newcommand{\Q}{\mathbb{Q}}
\newcommand{\R}{\mathbb{R}}
\newcommand{\C}{\mathbb{C}}
\newcommand{\F}{\mathbb{F}}          % 有限域

% 上半平面
\newcommand{\HH}{\mathbb{H}}         % 上半平面 H

% 模群与矩阵群
\newcommand{\SL}{\mathrm{SL}}
\newcommand{\GL}{\mathrm{GL}}
\newcommand{\PSL}{\mathrm{PSL}}
\newcommand{\PGL}{\mathrm{PGL}}
\newcommand{\Sp}{\mathrm{Sp}}
\newcommand{\Mat}{\mathrm{Mat}}

% 同余子群
\newcommand{\GammaO}{\Gamma_0}       % Γ₀(N)
\newcommand{\GammaI}{\Gamma_1}       % Γ₁(N)

% 模形式空间
\newcommand{\Mk}{\mathcal{M}}        % M_k
\newcommand{\Sk}{\mathcal{S}}        % S_k (尖点形式)

% Hecke 算子
\newcommand{\Tk}{T}                   % T_n

% L-函数
\newcommand{\Lfun}{L}

% 自守表示
\newcommand{\Aut}{\mathrm{Aut}}
\newcommand{\Gal}{\mathrm{Gal}}
\newcommand{\Frob}{\mathrm{Frob}}
\newcommand{\Ind}{\mathrm{Ind}}
\newcommand{\Hom}{\mathrm{Hom}}
\newcommand{\End}{\mathrm{End}}
\newcommand{\Tr}{\mathrm{Tr}}
\newcommand{\Nm}{\mathrm{Nm}}
\newcommand{\ord}{\mathrm{ord}}
\newcommand{\Res}{\mathrm{Res}}
\newcommand{\vol}{\mathrm{vol}}

% 其他
\newcommand{\diam}{\mathrm{diam}}
\newcommand{\Id}{\mathrm{Id}}
\DeclareMathOperator{\im}{im}
\DeclareMathOperator{\coker}{coker}
\DeclareMathOperator{\rank}{rank}
\DeclareMathOperator{\Div}{Div}
\DeclareMathOperator{\divisor}{div}
\DeclareMathOperator{\genus}{genus}


% ============================================================================
\begin{document}

\raggedbottom

% --- 封面 ---
\begin{titlepage}
  \begin{tikzpicture}[remember picture, overlay]
    % 顶部窄线
    \fill[black!60] (current page.north west) rectangle
      ([yshift=-0.6cm]current page.north east);
    % 深蓝标题色带
    \fill[blue!40!black] ([yshift=-2.5cm]current page.north west) rectangle
      ([yshift=-7.5cm]current page.north east);
    % 底部窄线
    \fill[black!60] ([yshift=1.8cm]current page.south west) rectangle
      ([yshift=2.2cm]current page.south east);
  \end{tikzpicture}

  \centering
  \vspace*{2.2cm}

  {\color{white}\fontsize{38}{46}\selectfont\bfseries 模形式与费马大定理\par}
  \vspace{0.5cm}
  {\color{white!90}\LARGE 从 $\SL_2(\Z)$ 到 Wiles 定理\par}

  \vspace{3cm}

  {\large\itshape 基于 Diamond \& Shurman 的学习笔记\par}
  \vspace{0.3cm}
  {\normalsize 力求每一个定理都给出详细证明\par}

  \vfill

  {\Large\bfseries SZY\par}
  \vspace{0.3cm}
  {\normalsize \today\par}
  \vspace{2.5cm}
\end{titlepage}

% --- 目录 ---
\tableofcontents
\newpage

% ============================================================================
% 第一部分：预备知识
% ============================================================================
\part{预备知识}

% ============================================================================
\chapter{复分析基础回顾}
\label{ch:complex-analysis}
% ============================================================================

本章回顾模形式理论所需的复分析背景。读者若已熟悉这些内容，可跳至
\cref{ch:elliptic-curves-intro}。

% ----------------------------------------------------------------------------
\section{全纯函数与亚纯函数}
% ----------------------------------------------------------------------------

\begin{definition}[全纯函数]
设 $U \subset \C$ 为开集，函数 $f\colon U \to \C$ 在 $z_0 \in U$ 处
\textbf{全纯}（holomorphic），若极限
\[
  f'(z_0) = \lim_{h \to 0} \frac{f(z_0 + h) - f(z_0)}{h}
\]
存在。若 $f$ 在 $U$ 上每一点都全纯，则称 $f$ 在 $U$ 上全纯。
\end{definition}

% ----------------------------------------------------------------------------
\section{Laurent 级数与留数}
% ----------------------------------------------------------------------------

% TODO: Laurent 展开、留数定理

% ----------------------------------------------------------------------------
\section{黎曼面初步}
% ----------------------------------------------------------------------------

% TODO: 黎曼面、紧化、亏格

% ----------------------------------------------------------------------------
\section{椭圆函数}
% ----------------------------------------------------------------------------

% TODO: 格、Weierstrass ℘ 函数、椭圆函数的周期性

% ----------------------------------------------------------------------------
\section{习题}
% ----------------------------------------------------------------------------

% ============================================================================
\chapter{椭圆曲线初步}
\label{ch:elliptic-curves-intro}
% ============================================================================

本章介绍椭圆曲线的基本概念，为后续将模形式与椭圆曲线联系起来做准备。

% ----------------------------------------------------------------------------
\section{Weierstrass 方程}
% ----------------------------------------------------------------------------

\begin{definition}[椭圆曲线]
域 $K$ 上的一条\textbf{椭圆曲线}是一条亏格为 $1$ 的光滑射影曲线，
并指定一个 $K$-有理点 $O$ 作为单位元。其 Weierstrass 方程为
\[
  y^2 = x^3 + ax + b, \quad a, b \in K, \quad 4a^3 + 27b^2 \neq 0.
\]
\end{definition}

% ----------------------------------------------------------------------------
\section{群结构}
% ----------------------------------------------------------------------------

% TODO: 椭圆曲线上的加法、弦切法

% ----------------------------------------------------------------------------
\section{同源（Isogeny）}
% ----------------------------------------------------------------------------

% TODO: 同源的定义、度、对偶同源

% ----------------------------------------------------------------------------
\section{椭圆曲线的 Tate 模}
% ----------------------------------------------------------------------------

% TODO: n-挠点、Tate 模 T_ℓ(E)

% ----------------------------------------------------------------------------
\section{$\Q$ 上椭圆曲线的约化}
% ----------------------------------------------------------------------------

% TODO: 好约化、坏约化、导子

% ----------------------------------------------------------------------------
\section{习题}
% ----------------------------------------------------------------------------


% ============================================================================
% 第二部分：$\SL_2(\Z)$ 上的模形式
% ============================================================================
\part{$\SL_2(\Z)$ 上的模形式}

% ============================================================================
\chapter{模群 $\SL_2(\Z)$}
\label{ch:modular-group}
% ============================================================================

% ----------------------------------------------------------------------------
\section{上半平面与模群的作用}
% ----------------------------------------------------------------------------

\begin{definition}[上半平面]
\textbf{上半平面}定义为
\[
  \HH = \{ \tau \in \C \mid \im(\tau) > 0 \}.
\]
\end{definition}

\begin{definition}[模群]
\textbf{模群} $\SL_2(\Z)$ 是所有行列式为 $1$ 的 $2 \times 2$ 整数矩阵组成的群：
\[
  \SL_2(\Z) = \left\{
    \begin{pmatrix} a & b \\ c & d \end{pmatrix}
    \;\middle|\; a, b, c, d \in \Z,\; ad - bc = 1
  \right\}.
\]
它通过\textbf{莫比乌斯变换}作用在 $\HH$ 上：
\[
  \gamma \cdot \tau = \frac{a\tau + b}{c\tau + d},
  \quad \gamma = \begin{pmatrix} a & b \\ c & d \end{pmatrix}.
\]
\end{definition}

% ----------------------------------------------------------------------------
\section{生成元 $S$ 和 $T$}
% ----------------------------------------------------------------------------

% TODO: S = ((0,-1),(1,0)), T = ((1,1),(0,1))，证明它们生成 SL₂(ℤ)

% ----------------------------------------------------------------------------
\section{基本域}
% ----------------------------------------------------------------------------

% TODO: 标准基本域 F 的描述与证明

% ----------------------------------------------------------------------------
\section{尖点}
% ----------------------------------------------------------------------------

% TODO: 扩展上半平面 H* = H ∪ Q ∪ {∞}，尖点的定义

% ----------------------------------------------------------------------------
\section{习题}
% ----------------------------------------------------------------------------

% ============================================================================
\chapter{$\SL_2(\Z)$ 的模形式}
\label{ch:modular-forms-sl2z}
% ============================================================================

% ----------------------------------------------------------------------------
\section{模形式的定义}
% ----------------------------------------------------------------------------

\begin{definition}[模形式]
设 $k$ 为非负偶整数。$\SL_2(\Z)$ 上权为 $k$ 的\textbf{模形式}是一个全纯函数
$f\colon \HH \to \C$，满足：
\begin{enumerate}
  \item \textbf{变换律}：对任意
    $\gamma = \begin{pmatrix} a & b \\ c & d \end{pmatrix} \in \SL_2(\Z)$，
    \[
      f\!\left(\frac{a\tau + b}{c\tau + d}\right) = (c\tau + d)^k f(\tau);
    \]
  \item \textbf{在尖点全纯}：$f$ 在 $\infty$ 处全纯，即 Fourier 展开
    \[
      f(\tau) = \sum_{n=0}^{\infty} a_n q^n, \quad q = e^{2\pi i \tau},
    \]
    的系数 $a_n = 0$ 对所有 $n < 0$ 成立。
\end{enumerate}
权 $k$ 的模形式空间记作 $\Mk_k(\SL_2(\Z))$。
\end{definition}

\begin{definition}[尖点形式]
若进一步 $a_0 = 0$（即 $f$ 在尖点处消失），则称 $f$ 为\textbf{尖点形式}。
尖点形式空间记作 $\Sk_k(\SL_2(\Z))$。
\end{definition}

% ----------------------------------------------------------------------------
\section{$q$-展开原理}
% ----------------------------------------------------------------------------

% TODO: q-展开原理的陈述与证明

% ----------------------------------------------------------------------------
\section{模形式的例子}
% ----------------------------------------------------------------------------

% TODO: Eisenstein 级数 G_k，判别形式 Δ，j-不变量

% ----------------------------------------------------------------------------
\section{赋值公式（Valence Formula）}
% ----------------------------------------------------------------------------

% TODO: 赋值公式的陈述与详细证明

% ----------------------------------------------------------------------------
\section{习题}
% ----------------------------------------------------------------------------

% ============================================================================
\chapter{Eisenstein 级数}
\label{ch:eisenstein-series}
% ============================================================================

% ----------------------------------------------------------------------------
\section{格与 Eisenstein 级数}
% ----------------------------------------------------------------------------

\begin{definition}[Eisenstein 级数]
对偶整数 $k \geq 4$，\textbf{Eisenstein 级数}定义为
\[
  G_k(\tau) = \sum_{\substack{(c,d) \in \Z^2 \\ (c,d) \neq (0,0)}}
  \frac{1}{(c\tau + d)^k}.
\]
\end{definition}

% TODO: 证明 G_k 是权 k 的模形式

% ----------------------------------------------------------------------------
\section{规范化 Eisenstein 级数 $E_k$}
% ----------------------------------------------------------------------------

% TODO: E_k = G_k / (2ζ(k))，Fourier 展开，Bernoulli 数

% ----------------------------------------------------------------------------
\section{判别形式 $\Delta$ 与 $j$-不变量}
% ----------------------------------------------------------------------------

% TODO: Δ = (E_4³ - E_6²)/1728，Ramanujan tau 函数，j = E_4³/Δ

% ----------------------------------------------------------------------------
\section{模形式空间的分级代数结构}
% ----------------------------------------------------------------------------

% TODO: M_*(SL₂(ℤ)) = C[E_4, E_6] 的证明

% ----------------------------------------------------------------------------
\section{习题}
% ----------------------------------------------------------------------------

% ============================================================================
\chapter{维数公式}
\label{ch:dimension-formula}
% ============================================================================

% ----------------------------------------------------------------------------
\section{$\Mk_k(\SL_2(\Z))$ 的维数}
% ----------------------------------------------------------------------------

\begin{theorem}[维数公式]
\label{thm:dim-formula-sl2z}
对非负偶整数 $k$，
\[
  \dim \Mk_k(\SL_2(\Z)) =
  \begin{cases}
    \lfloor k/12 \rfloor     & \text{若 } k \equiv 2 \pmod{12}, \\
    \lfloor k/12 \rfloor + 1 & \text{否则},
  \end{cases}
\]
且 $\dim \Sk_k(\SL_2(\Z)) = \dim \Mk_k(\SL_2(\Z)) - 1$（$k \geq 2$）。
\end{theorem}

% TODO: 详细证明（利用赋值公式）

% ----------------------------------------------------------------------------
\section{Riemann-Roch 定理的应用}
% ----------------------------------------------------------------------------

% TODO: 从 Riemann-Roch 出发推导维数公式

% ----------------------------------------------------------------------------
\section{同余子群的维数公式}
% ----------------------------------------------------------------------------

% TODO: Γ₀(N), Γ₁(N) 上模形式空间的维数（预告后续章节）

% ----------------------------------------------------------------------------
\section{习题}
% ----------------------------------------------------------------------------


% ============================================================================
% 第三部分：同余子群上的模形式
% ============================================================================
\part{同余子群上的模形式}

% ============================================================================
\chapter{同余子群}
\label{ch:congruence-subgroups}
% ============================================================================

% ----------------------------------------------------------------------------
\section{$\GammaO(N)$、$\GammaI(N)$ 和 $\Gamma(N)$ 的定义}
% ----------------------------------------------------------------------------

\begin{definition}[同余子群]
设 $N$ 为正整数。定义以下 $\SL_2(\Z)$ 的子群：
\begin{align}
  \Gamma(N) &= \left\{
    \begin{pmatrix} a & b \\ c & d \end{pmatrix} \in \SL_2(\Z)
    \;\middle|\;
    \begin{pmatrix} a & b \\ c & d \end{pmatrix}
    \equiv \begin{pmatrix} 1 & 0 \\ 0 & 1 \end{pmatrix} \pmod{N}
  \right\}, \\[6pt]
  \GammaI(N) &= \left\{
    \begin{pmatrix} a & b \\ c & d \end{pmatrix} \in \SL_2(\Z)
    \;\middle|\;
    \begin{pmatrix} a & b \\ c & d \end{pmatrix}
    \equiv \begin{pmatrix} 1 & * \\ 0 & 1 \end{pmatrix} \pmod{N}
  \right\}, \\[6pt]
  \GammaO(N) &= \left\{
    \begin{pmatrix} a & b \\ c & d \end{pmatrix} \in \SL_2(\Z)
    \;\middle|\;
    c \equiv 0 \pmod{N}
  \right\}.
\end{align}
包含关系为 $\Gamma(N) \trianglelefteq \GammaI(N) \trianglelefteq \GammaO(N)
\leq \SL_2(\Z)$。
\end{definition}

% ----------------------------------------------------------------------------
\section{指标与陪集分解}
% ----------------------------------------------------------------------------

% TODO: 各子群在 SL₂(ℤ) 中的指标公式

% ----------------------------------------------------------------------------
\section{基本域与尖点}
% ----------------------------------------------------------------------------

% TODO: 同余子群的基本域、尖点的分类

% ----------------------------------------------------------------------------
\section{习题}
% ----------------------------------------------------------------------------

% ============================================================================
\chapter{同余子群上的模形式}
\label{ch:modular-forms-congruence}
% ============================================================================

% ----------------------------------------------------------------------------
\section{同余子群上的模形式与尖点形式}
% ----------------------------------------------------------------------------

\begin{definition}
设 $\Gamma$ 为 $\SL_2(\Z)$ 的一个同余子群。权 $k$ 的\textbf{模形式}
$f \in \Mk_k(\Gamma)$ 是全纯函数 $f\colon \HH \to \C$，满足：
\begin{enumerate}
  \item 对任意 $\gamma = \begin{pmatrix} a & b \\ c & d \end{pmatrix}
    \in \Gamma$，有 $f(\gamma \cdot \tau) = (c\tau + d)^k f(\tau)$；
  \item $f$ 在 $\Gamma$ 的每个尖点处全纯。
\end{enumerate}
\end{definition}

% ----------------------------------------------------------------------------
\section{$\Mk_k(\GammaO(N))$ 与 $\Mk_k(\GammaI(N))$ 的维数公式}
% ----------------------------------------------------------------------------

% TODO: 一般同余子群的维数公式（Riemann-Roch 方法）

% ----------------------------------------------------------------------------
\section{Diamond 算子与 $\langle d \rangle$ 作用}
% ----------------------------------------------------------------------------

% TODO: Diamond 算子的定义与性质

% ----------------------------------------------------------------------------
\section{oldforms 与 newforms（预告）}
% ----------------------------------------------------------------------------

% TODO: 简要预告 Atkin-Lehner 理论

% ----------------------------------------------------------------------------
\section{习题}
% ----------------------------------------------------------------------------

% ============================================================================
\chapter{Dirichlet 特征与 Nebentypus}
\label{ch:dirichlet-characters}
% ============================================================================

% ----------------------------------------------------------------------------
\section{Dirichlet 特征}
% ----------------------------------------------------------------------------

\begin{definition}[Dirichlet 特征]
模 $N$ 的 \textbf{Dirichlet 特征}是群同态
$\chi\colon (\Z/N\Z)^\times \to \C^\times$。
将其扩展到 $\Z$ 上：若 $\gcd(n, N) > 1$，令 $\chi(n) = 0$。
\end{definition}

% TODO: 原特征、导子、Gauss 和

% ----------------------------------------------------------------------------
\section{带 Nebentypus 的模形式}
% ----------------------------------------------------------------------------

% TODO: M_k(Γ₀(N), χ) 的定义、与 M_k(Γ₁(N)) 的分解

% ----------------------------------------------------------------------------
\section{扭（Twist）}
% ----------------------------------------------------------------------------

% TODO: 模形式的 Dirichlet 特征扭

% ----------------------------------------------------------------------------
\section{习题}
% ----------------------------------------------------------------------------


% ============================================================================
% 第四部分：Hecke 理论
% ============================================================================
\part{Hecke 理论}

% ============================================================================
\chapter{Hecke 算子}
\label{ch:hecke-operators}
% ============================================================================

% ----------------------------------------------------------------------------
\section{双陪集与 Hecke 算子的定义}
% ----------------------------------------------------------------------------

\begin{definition}[Hecke 算子]
对素数 $p$，作用在 $\Mk_k(\SL_2(\Z))$ 上的 \textbf{Hecke 算子} $T_p$ 定义为
\[
  (T_p f)(\tau) = p^{k-1} f(p\tau)
  + \frac{1}{p} \sum_{j=0}^{p-1} f\!\left(\frac{\tau + j}{p}\right).
\]
等价地，用 $q$-展开：若 $f = \sum a_n q^n$，则
\[
  T_p f = \sum_{n=0}^{\infty} \bigl( a_{np} + p^{k-1} a_{n/p} \bigr) q^n,
\]
其中约定 $n/p \notin \Z$ 时 $a_{n/p} = 0$。
\end{definition}

% ----------------------------------------------------------------------------
\section{Hecke 算子的性质}
% ----------------------------------------------------------------------------

% TODO: T_m T_n = T_{mn} 当 gcd(m,n)=1，Hecke 代数的交换性

% ----------------------------------------------------------------------------
\section{Hecke 本征形式}
% ----------------------------------------------------------------------------

% TODO: 本征形式的定义、a₁ = 1 的规范化、Fourier 系数的乘性

% ----------------------------------------------------------------------------
\section{同余子群上的 Hecke 算子}
% ----------------------------------------------------------------------------

% TODO: T_n 在 Γ₀(N) 上的定义，Diamond 算子 ⟨d⟩ 与 T_p 的关系

% ----------------------------------------------------------------------------
\section{习题}
% ----------------------------------------------------------------------------

% ============================================================================
\chapter{Petersson 内积与谱分解}
\label{ch:petersson-inner-product}
% ============================================================================

% ----------------------------------------------------------------------------
\section{Petersson 内积}
% ----------------------------------------------------------------------------

\begin{definition}[Petersson 内积]
对 $f, g \in \Sk_k(\Gamma)$（至少一个为尖点形式），定义
\[
  \langle f, g \rangle = \frac{1}{[\SL_2(\Z):\Gamma]}
  \int_{\Gamma \backslash \HH} f(\tau) \overline{g(\tau)}\, (\im \tau)^k
  \, \frac{dx\,dy}{y^2},
\]
其中 $\tau = x + iy$。
\end{definition}

% TODO: 证明内积的收敛性

% ----------------------------------------------------------------------------
\section{Hecke 算子的自伴性}
% ----------------------------------------------------------------------------

% TODO: 证明 T_n 关于 Petersson 内积是自伴算子（当 gcd(n,N)=1）

% ----------------------------------------------------------------------------
\section{谱定理与 Hecke 本征形式基}
% ----------------------------------------------------------------------------

% TODO: S_k 有 Hecke 本征形式组成的正交基

% ----------------------------------------------------------------------------
\section{习题}
% ----------------------------------------------------------------------------

% ============================================================================
\chapter{Atkin-Lehner 理论与新形式}
\label{ch:newforms}
% ============================================================================

% ----------------------------------------------------------------------------
\section{旧形式（Oldforms）}
% ----------------------------------------------------------------------------

% TODO: 从低 level 通过 d|N 映射得到的模形式

% ----------------------------------------------------------------------------
\section{新形式（Newforms）}
% ----------------------------------------------------------------------------

\begin{definition}[新形式]
$\Sk_k(\GammaO(N))$ 中的\textbf{新形式}（newform）是旧形式空间的正交补中的
规范化 Hecke 本征形式。新形式空间记为 $\Sk_k^{\mathrm{new}}(\GammaO(N))$。
\end{definition}

% TODO: 新形式的存在性与唯一性（Multiplicity One 定理）

% ----------------------------------------------------------------------------
\section{Atkin-Lehner 对合}
% ----------------------------------------------------------------------------

% TODO: W_Q 算子、函数方程中的符号

% ----------------------------------------------------------------------------
\section{Strong Multiplicity One 定理}
% ----------------------------------------------------------------------------

% TODO: 陈述与证明

% ----------------------------------------------------------------------------
\section{习题}
% ----------------------------------------------------------------------------


% ============================================================================
% 第五部分：$L$-函数与 Theta 函数
% ============================================================================
\part{$L$-函数与 Theta 函数}

% ============================================================================
\chapter{模形式的 $L$-函数}
\label{ch:l-functions}
% ============================================================================

% ----------------------------------------------------------------------------
\section{$L$-函数的定义}
% ----------------------------------------------------------------------------

\begin{definition}
设 $f(\tau) = \sum_{n=1}^{\infty} a_n q^n \in \Sk_k(\GammaO(N))$ 为规范化新形式，
其 \textbf{$L$-函数}定义为 Dirichlet 级数
\[
  L(f, s) = \sum_{n=1}^{\infty} \frac{a_n}{n^s},
  \quad \Re(s) > \frac{k+1}{2}.
\]
\end{definition}

% ----------------------------------------------------------------------------
\section{Euler 乘积}
% ----------------------------------------------------------------------------

% TODO: L(f,s) = ∏_p (1 - a_p p^{-s} + p^{k-1-2s})^{-1}

% ----------------------------------------------------------------------------
\section{Mellin 变换}
% ----------------------------------------------------------------------------

% TODO: L(f,s) 与 f 的 Mellin 变换的关系

% ----------------------------------------------------------------------------
\section{Hecke 的逆定理}
% ----------------------------------------------------------------------------

% TODO: Hecke 逆定理的陈述（从 L-函数恢复模形式）

% ----------------------------------------------------------------------------
\section{习题}
% ----------------------------------------------------------------------------

% ============================================================================
\chapter{函数方程与解析延拓}
\label{ch:functional-equation}
% ============================================================================

% ----------------------------------------------------------------------------
\section{完备 $L$-函数 $\Lambda(f, s)$}
% ----------------------------------------------------------------------------

% TODO: Γ 因子、完备化 Λ(f,s) = N^{s/2} (2π)^{-s} Γ(s) L(f,s)

% ----------------------------------------------------------------------------
\section{函数方程}
% ----------------------------------------------------------------------------

\begin{theorem}[函数方程]
设 $f \in \Sk_k^{\mathrm{new}}(\GammaO(N))$ 为规范化新形式，则完备
$L$-函数 $\Lambda(f, s)$ 可以解析延拓到整个复平面，且满足函数方程
\[
  \Lambda(f, s) = \epsilon \cdot \Lambda(f, k - s),
\]
其中 $\epsilon = \pm 1$ 为 $f$ 的符号（由 Atkin-Lehner 对合决定）。
\end{theorem}

% TODO: 详细证明

% ----------------------------------------------------------------------------
\section{$L(f, s)$ 在 $s = k/2$ 处的特殊值}
% ----------------------------------------------------------------------------

% TODO: BSD 猜想的预告

% ----------------------------------------------------------------------------
\section{习题}
% ----------------------------------------------------------------------------

% ============================================================================
\chapter{Theta 函数}
\label{ch:theta-functions}
% ============================================================================

% ----------------------------------------------------------------------------
\section{Jacobi Theta 函数}
% ----------------------------------------------------------------------------

\begin{definition}
\textbf{Jacobi theta 函数}定义为
\[
  \theta(\tau) = \sum_{n \in \Z} q^{n^2} = 1 + 2\sum_{n=1}^{\infty} q^{n^2},
  \quad q = e^{\pi i \tau}.
\]
\end{definition}

% TODO: 证明 θ 是 Γ₀(4) 上权 1/2 的模形式

% ----------------------------------------------------------------------------
\section{二次型的表示数}
% ----------------------------------------------------------------------------

% TODO: r_k(n) 与 θ 函数的幂次

% ----------------------------------------------------------------------------
\section{半整数权模形式}
% ----------------------------------------------------------------------------

% TODO: 半整数权模形式的定义、Shimura 对应

% ----------------------------------------------------------------------------
\section{Theta 级数与格}
% ----------------------------------------------------------------------------

% TODO: 格的 theta 级数、偶幺模格

% ----------------------------------------------------------------------------
\section{习题}
% ----------------------------------------------------------------------------


% ============================================================================
% 第六部分：模曲线
% ============================================================================
\part{模曲线}

% ============================================================================
\chapter{模曲线作为黎曼面}
\label{ch:modular-curves-riemann}
% ============================================================================

% ----------------------------------------------------------------------------
\section{商空间 $\Gamma \backslash \HH^*$}
% ----------------------------------------------------------------------------

% TODO: H* = H ∪ Q ∪ {∞}，商空间的拓扑、复结构

% ----------------------------------------------------------------------------
\section{紧黎曼面 $X(\Gamma)$}
% ----------------------------------------------------------------------------

\begin{definition}
设 $\Gamma$ 为 $\SL_2(\Z)$ 的同余子群，\textbf{模曲线} $X(\Gamma)$ 定义为
紧黎曼面 $\Gamma \backslash \HH^*$。特别地：
\begin{itemize}
  \item $X_0(N) = X(\GammaO(N))$,
  \item $X_1(N) = X(\GammaI(N))$,
  \item $X(N) = X(\Gamma(N))$.
\end{itemize}
\end{definition}

% ----------------------------------------------------------------------------
\section{亏格公式}
% ----------------------------------------------------------------------------

% TODO: 用 Riemann-Hurwitz 计算 X₀(N) 的亏格

% ----------------------------------------------------------------------------
\section{模形式与微分形式}
% ----------------------------------------------------------------------------

% TODO: S_k(Γ) ≅ H⁰(X(Γ), Ω^{⊗(k/2)}) 的同构

% ----------------------------------------------------------------------------
\section{习题}
% ----------------------------------------------------------------------------

% ============================================================================
\chapter{模曲线作为代数曲线}
\label{ch:modular-curves-algebraic}
% ============================================================================

% ----------------------------------------------------------------------------
\section{GAGA 原理}
% ----------------------------------------------------------------------------

% TODO: 紧黎曼面 ↔ 光滑射影代数曲线

% ----------------------------------------------------------------------------
\section{$X_0(N)$ 与 $X_1(N)$ 的代数模型}
% ----------------------------------------------------------------------------

% TODO: 作为 Q 上的代数曲线、有理点

% ----------------------------------------------------------------------------
\section{模曲线上的 Hecke 对应}
% ----------------------------------------------------------------------------

% TODO: Hecke 对应的代数定义、与 Hecke 算子的关系

% ----------------------------------------------------------------------------
\section{Jacobian 与 Abel 簇}
% ----------------------------------------------------------------------------

% TODO: J₀(N) = Jac(X₀(N))

% ----------------------------------------------------------------------------
\section{习题}
% ----------------------------------------------------------------------------

% ============================================================================
\chapter{模空间诠释}
\label{ch:moduli-interpretation}
% ============================================================================

% ----------------------------------------------------------------------------
\section{$Y_1(N)$ 作为椭圆曲线的模空间}
% ----------------------------------------------------------------------------

% TODO: Y₁(N) 参数化 (E, P)，其中 P 是 N-挠点

% ----------------------------------------------------------------------------
\section{$Y_0(N)$ 与循环同源}
% ----------------------------------------------------------------------------

% TODO: Y₀(N) 参数化 (E → E')，核为 Z/NZ 的同源

% ----------------------------------------------------------------------------
\section{模问题的层论观点}
% ----------------------------------------------------------------------------

% TODO: 简介表示函子（representable functor）

% ----------------------------------------------------------------------------
\section{Hecke 对应的模诠释}
% ----------------------------------------------------------------------------

% TODO: T_p 对应 p-同源的选取

% ----------------------------------------------------------------------------
\section{习题}
% ----------------------------------------------------------------------------


% ============================================================================
% 第七部分：椭圆曲线与 Galois 表示
% ============================================================================
\part{椭圆曲线与 Galois 表示}

% ============================================================================
\chapter{椭圆曲线的算术}
\label{ch:elliptic-curve-arithmetic}
% ============================================================================

% ----------------------------------------------------------------------------
\section{$\Q$ 上椭圆曲线的 Mordell-Weil 定理}
% ----------------------------------------------------------------------------

\begin{theorem}[Mordell-Weil]
设 $E/\Q$ 为椭圆曲线，则 $E(\Q)$ 是有限生成 Abel 群：
\[
  E(\Q) \cong \Z^r \oplus E(\Q)_{\mathrm{tors}},
\]
其中 $r = \rank(E/\Q) \geq 0$。
\end{theorem}

% TODO: 证明思路概述（高度函数、弱 Mordell-Weil）

% ----------------------------------------------------------------------------
\section{约化类型与导子}
% ----------------------------------------------------------------------------

% TODO: 好约化、乘性约化、加性约化；导子 N_E 的定义

% ----------------------------------------------------------------------------
\section{椭圆曲线的 $L$-函数}
% ----------------------------------------------------------------------------

% TODO: L(E, s) = ∏ L_p(E, s)，Hasse-Weil L-函数

% ----------------------------------------------------------------------------
\section{BSD 猜想}
% ----------------------------------------------------------------------------

% TODO: 陈述 BSD 猜想、已知结果

% ----------------------------------------------------------------------------
\section{习题}
% ----------------------------------------------------------------------------

% ============================================================================
\chapter{Galois 表示}
\label{ch:galois-representations}
% ============================================================================

% ----------------------------------------------------------------------------
\section{绝对 Galois 群 $G_\Q$}
% ----------------------------------------------------------------------------

% TODO: G_Q = Gal(Q̄/Q)，分解群、惯性群、Frobenius 元素

% ----------------------------------------------------------------------------
\section{$\ell$-进 Galois 表示}
% ----------------------------------------------------------------------------

\begin{definition}
一个 \textbf{$\ell$-进 Galois 表示}是连续群同态
\[
  \rho\colon G_\Q \to \GL_n(\Q_\ell).
\]
\end{definition}

% TODO: 非分歧表示、Artin 表示、半简单化

% ----------------------------------------------------------------------------
\section{椭圆曲线的 Galois 表示}
% ----------------------------------------------------------------------------

% TODO: ρ_{E,ℓ}: G_Q → GL₂(Q_ℓ) 来自 Tate 模

% TODO: 在好约化素数 p 处：Tr(Frob_p) = a_p，det(Frob_p) = p

% ----------------------------------------------------------------------------
\section{模形式的 Galois 表示}
% ----------------------------------------------------------------------------

% TODO: Deligne 定理：权 k 新形式 f 给出 ρ_f: G_Q → GL₂(Q̄_ℓ)

% ----------------------------------------------------------------------------
\section{习题}
% ----------------------------------------------------------------------------

% ============================================================================
\chapter{Eichler-Shimura 关系}
\label{ch:eichler-shimura}
% ============================================================================

% ----------------------------------------------------------------------------
\section{模曲线的好约化}
% ----------------------------------------------------------------------------

% TODO: X₀(N) 在 p ∤ N 时的好约化

% ----------------------------------------------------------------------------
\section{Eichler-Shimura 关系}
% ----------------------------------------------------------------------------

\begin{theorem}[Eichler-Shimura]
设 $p \nmid N$ 为素数，则在 $X_0(N)$ 的模 $p$ 约化上，
Hecke 对应 $T_p$ 满足
\[
  T_p = \Frob_p + \frac{p}{\Frob_p}
\]
（作为 $J_0(N)$ 的 $\F_p$-自同态，这里 $\Frob_p$ 是 Frobenius 自同态）。
\end{theorem}

% TODO: 详细证明

% ----------------------------------------------------------------------------
\section{推论：$a_p$ 与 Frobenius 迹}
% ----------------------------------------------------------------------------

% TODO: 由 Eichler-Shimura 得 a_p(f) = Tr(Frob_p | V_ℓ(A_f))

% ----------------------------------------------------------------------------
\section{与 Weil 猜想的联系}
% ----------------------------------------------------------------------------

% TODO: Hasse 界 |a_p| ≤ 2√p（Ramanujan-Petersson 猜想）

% ----------------------------------------------------------------------------
\section{习题}
% ----------------------------------------------------------------------------


% ============================================================================
% 第八部分：模性定理与费马大定理
% ============================================================================
\part{模性定理与费马大定理}

% ============================================================================
\chapter{模性猜想}
\label{ch:modularity-conjecture}
% ============================================================================

% ----------------------------------------------------------------------------
\section{Taniyama-Shimura-Weil 猜想的陈述}
% ----------------------------------------------------------------------------

\begin{conjecture}[模性猜想 / Taniyama-Shimura-Weil]
\label{conj:modularity}
设 $E/\Q$ 为椭圆曲线，导子为 $N$。则存在权 $2$ 的规范化新形式
$f \in \Sk_2^{\mathrm{new}}(\GammaO(N))$，使得
\[
  L(E, s) = L(f, s).
\]
等价地，存在非常值态射 $\varphi\colon X_0(N) \to E$（定义在 $\Q$ 上）。
\end{conjecture}

% ----------------------------------------------------------------------------
\section{猜想的等价形式}
% ----------------------------------------------------------------------------

% TODO: Galois 表示版本：ρ_{E,ℓ} ≅ ρ_{f,ℓ}
% TODO: 同源版本：E 是 J₀(N) 的商

% ----------------------------------------------------------------------------
\section{已知的特殊情形}
% ----------------------------------------------------------------------------

% TODO: CM 椭圆曲线（Hecke-Deuring-Shimura）

% ----------------------------------------------------------------------------
\section{从模性猜想到费马大定理}
% ----------------------------------------------------------------------------

% TODO: Frey 曲线、Ribet 定理（level lowering）、Serre 猜想的 ε-情形

% ----------------------------------------------------------------------------
\section{习题}
% ----------------------------------------------------------------------------

% ============================================================================
\chapter{Wiles 的证明策略}
\label{ch:wiles-proof}
% ============================================================================

% ----------------------------------------------------------------------------
\section{证明的整体架构}
% ----------------------------------------------------------------------------

% TODO: 
% Wiles 证明模性猜想（半稳定情形）的路线图：
% 1. 将问题归结为 ρ̄_{E,3} 的模性提升
% 2. R = T 定理
% 3. Taylor-Wiles 方法

% ----------------------------------------------------------------------------
\section{残余表示与模性提升}
% ----------------------------------------------------------------------------

% TODO: 
% - 残余表示 ρ̄: G_Q → GL₂(F_ℓ)
% - Serre 猜想的背景
% - 从 ρ̄ 的模性到 ρ 的模性（提升问题）

% ----------------------------------------------------------------------------
\section{变形环与 Hecke 代数}
% ----------------------------------------------------------------------------

% TODO:
% - 万有变形环 R
% - Hecke 代数 T
% - R = T 定理的陈述

% ----------------------------------------------------------------------------
\section{Taylor-Wiles 方法}
% ----------------------------------------------------------------------------

% TODO:
% - 辅助素数集（Taylor-Wiles primes）
% - 极限论证
% - R = T 的证明思路

% ----------------------------------------------------------------------------
\section{Langlands-Tunnell 定理与 3-5 互换}
% ----------------------------------------------------------------------------

% TODO:
% - 用 Langlands-Tunnell 启动 ℓ = 3 的情形
% - 当 ρ̄_{E,3} 不适用时切换到 ℓ = 5

% ----------------------------------------------------------------------------
\section{习题}
% ----------------------------------------------------------------------------

% ============================================================================
\chapter{费马大定理}
\label{ch:flt}
% ============================================================================

% ----------------------------------------------------------------------------
\section{费马大定理的陈述}
% ----------------------------------------------------------------------------

\begin{theorem}[费马大定理, Wiles 1995]
\label{thm:flt}
方程
\[
  x^n + y^n = z^n
\]
在 $n \geq 3$ 时没有正整数解。
\end{theorem}

% ----------------------------------------------------------------------------
\section{历史背景}
% ----------------------------------------------------------------------------

% TODO:
% - Fermat 的批注（1637）
% - 特殊情形：n=3 (Euler), n=4 (Fermat), n=5 (Dirichlet-Legendre), 
%   n=7 (Lamé), 正则素数 (Kummer)
% - 计算机验证的进展

% ----------------------------------------------------------------------------
\section{Frey 曲线}
% ----------------------------------------------------------------------------

% TODO:
% - 若 aᵖ + bᵖ = cᵖ，构造 Frey 曲线 E: y² = x(x-aᵖ)(x+bᵖ)
% - Frey 曲线的性质：半稳定，导子 N = rad(abc)²

% ----------------------------------------------------------------------------
\section{Ribet 定理（Level Lowering）}
% ----------------------------------------------------------------------------

\begin{theorem}[Ribet, 1990]
设 $E/\Q$ 为半稳定椭圆曲线。若 $E$ 模性（即存在 $f \in \Sk_2(\GammaO(N_E))$
使得 $\rho_{E,p} \cong \rho_{f,p}$），则可以降低 level：
存在 $g \in \Sk_2(\GammaO(2))$ 使得 $\rho_{E,p} \cong \rho_{g,p}$。
\end{theorem}

% TODO: 
% - 但 S₂(Γ₀(2)) = 0，矛盾！
% - 因此 Frey 曲线不存在 → FLT 成立

% ----------------------------------------------------------------------------
\section{完整的证明逻辑链}
% ----------------------------------------------------------------------------

% TODO:
% 假设 aᵖ + bᵖ = cᵖ 有正整数解（p ≥ 5 为奇素数）
% → 构造 Frey 曲线 E
% → E 半稳定（Frey）
% → E 模性（Wiles-Taylor 的模性定理）
% → Level lowering 到 Γ₀(2)（Ribet 定理）
% → S₂(Γ₀(2)) = 0，矛盾
% → FLT 成立

% ----------------------------------------------------------------------------
\section{后续发展}
% ----------------------------------------------------------------------------

% TODO:
% - 完全模性定理（Breuil-Conrad-Diamond-Taylor, 2001）
% - Serre 猜想（Khare-Wintenberger, 2009）
% - 推广到全实域（Freitas-Le Hung-Siksek, 2015）

% ----------------------------------------------------------------------------
\section{习题}
% ----------------------------------------------------------------------------


% ============================================================================
% 附录
% ============================================================================
\appendix
\part{附录}

\include{chapters/app-a-algebra}
% ============================================================================
\chapter{表示论基础}
\label{app:representation-theory}
% ============================================================================

% TODO:
% - 群表示的基本概念：不可约、完全可约、Schur 引理
% - 特征理论
% - ℓ-进表示与 Galois 表示
% - 自守表示简介（GL₂ 上的光滑表示、容许表示）

\include{chapters/app-c-history}

% --- 参考文献 ---
\bibliographystyle{plain}
% \bibliography{refs}  % 待添加 refs.bib

% --- 索引 ---
\printindex

\end{document}
